\documentclass[11pt,letterpaper]{article}

\usepackage[margin=1.1in]{geometry}
\usepackage{amsmath, amssymb}
\usepackage{mathtools}
\usepackage{bm}
\usepackage{booktabs}
\usepackage{array}
\usepackage{enumitem}
\usepackage{xcolor}
\usepackage{tcolorbox}
\usepackage{hyperref}
\usepackage[protrusion=true,expansion=false]{microtype}
\usepackage{titlesec}
\usepackage{fancyhdr}
\usepackage{longtable}
\usepackage{graphicx}
\usepackage{float}
\graphicspath{{figs/}{SeanEarl_DiscFlight/P5_Synthesis/figs/}}

\definecolor{accent}{RGB}{30,80,60}
\definecolor{accent2}{RGB}{20,100,100}
\definecolor{boxbg}{RGB}{240,238,232}

\hypersetup{
  colorlinks=true,
  linkcolor=accent,
  citecolor=accent,
  urlcolor=accent2,
  pdftitle={Disc Flight in a Lakefront Internal Boundary Layer (Draft)},
  pdfauthor={charlie}
}

\titleformat{\section}
  {\large\bfseries\color{accent}}
  {\thesection.}{0.6em}{}[\vspace{-2pt}\rule{\linewidth}{0.6pt}]
\titleformat{\subsection}
  {\normalsize\bfseries\color{accent!80!black}}
  {\thesubsection.}{0.5em}{}
\titleformat{\subsubsection}
  {\normalsize\itshape}
  {\thesubsubsection.}{0.4em}{}

\pagestyle{fancy}
\fancyhf{}
\rhead{\small\color{gray}Project 1 $\cdot$ Frisbee + Loyola Wind}
\lhead{\small\color{gray}Manuscript Draft v0.1}
\cfoot{\thepage}
\renewcommand{\headrulewidth}{0.3pt}

\setlength{\parskip}{0.5\baselineskip}
\setlength{\parindent}{0pt}

\title{\bfseries Disc Flight in a Lakefront Internal Boundary Layer:\\Coupling Analytical Frisbee Aerodynamics with Site-Specific Urban CFD\\over Loyola University's Sean Earl Field}
\author{charlie\\\small affiliation TBD}
\date{Draft v0.1 \ $\cdot$\ \today}

\begin{document}
\maketitle

\noindent\textbf{Companion document:} \emph{Project1\_AnalyticalTreatment.tex} (rigorous 6-DOF derivation, stability, gyroscopic precession).

\section*{How to read this document}

Three layers stacked: \textbf{feasibility memo} (\S\ref{sec:feas}), \textbf{project plan} (\S\ref{sec:plan}), and \textbf{paper first draft} (\S\ref{sec:abstract} onward). Read \S\ref{sec:feas} first; if you disagree with the scope changes proposed there, the draft from \S\ref{sec:abstract} onward needs to be reframed.

\begin{tcolorbox}[colback=boxbg,colframe=accent,boxrule=0.6pt,left=6pt,
                  right=6pt,top=4pt,bottom=4pt]
\small\textbf{Status of this version.} \S\ref{sec:feas} and \S\ref{sec:plan}
are the \emph{original planning memo} and are retained unedited as a record of
intent. Where the executed work departs from them --- and it does, in the
number of CFD directions, the solver version, the building heights and the site
description --- \textbf{the paper sections from \S\ref{sec:abstract} onward are
authoritative}, and \S\ref{sec:limits} states the delivered scope explicitly.
All numbers quoted from \S\ref{sec:abstract} onward come from a converged run
or a passing test; nothing is projected.
\end{tcolorbox}

\section{Feasibility check}
\label{sec:feas}

\subsection{What the notes propose}

\begin{quote}
(1) Analytical treatment of ideal conditions for throwing a frisbee.\\
(2) Build a realistic 3D map of the terrain surrounding Loyola field --- houses, buildings --- simulate wind off the lake spanning 180$^\circ$, generate 3D flow fields across the turf to capture vortices.\\
\emph{Goal:} determine the best conditions for this constrained terrain that best allow for frisbee.\\
\emph{Tools:} OpenFOAM or COMSOL.
\end{quote}

\subsection{Verdict}

\begin{center}
\footnotesize
\begin{tabular}{@{}p{0.32\linewidth}p{0.15\linewidth}p{0.48\linewidth}@{}}
\toprule
\textbf{Component} & \textbf{Verdict} & \textbf{Why} \\
\midrule
Analytical 6-DOF disc model & Solid & Hummel (2003), Crowther \& Potts (2007) give a complete validated 6-DOF formulation with measured coefficients. \\
RANS urban CFD, 13 wind directions $\times$ 15$^\circ$ & Tractable & $\sim$3--10 M cells, 2--8 h per direction on 32 cores, total 1--4 days. Standard COST 732 / AIJ practice; OpenFOAM \texttt{simpleFoam} with k-$\omega$ SST. \\
LES of the same domain & Not realistic for full sweep & $\Delta x \sim 1$--$2$ m near buildings $\Rightarrow$ 50--200 M cells per case. Single direction is feasible (weeks). \\
Coupling CFD wind field into disc trajectory & Novel and tractable & Not done in published literature; the Lukes et al.\ (2022) disc-golf code uses uniform freestream. \emph{This is the novelty hook.} \\
Throwability metric over the field & Tractable as descriptive optimization & Not real-time forecasting; compute mean velocity + TKE at disc altitude, then evaluate trajectory robustness for canonical throws. \\
\bottomrule
\end{tabular}
\end{center}

\subsection{Items in your notes that need adjustment}

\begin{itemize}[leftmargin=*,itemsep=2pt]
  \item \textbf{``Loyola field''} --- your phrasing (``off the lake'') uniquely matches \emph{Sean Earl Field} on the Lake Shore Campus, $\sim$150--250 m from Lake Michigan, not Hoyne Field (a mile inland). \textbf{Confirmed by you: Sean Earl.}
  \item \textbf{``180$^\circ$ span''} --- Sean Earl's lake-facing side is roughly east; a meaningful ``off-lake'' sweep is NNE through ESE ($\sim$90$^\circ$ centered on east). A full 180$^\circ$ is fine if you also characterize the prevailing westerly (off-land) cases for contrast --- think of it as two regimes, not one continuous sweep.
  \item \textbf{``Range finder'' to build the 3D map} --- you don't need a personal LiDAR pass. USGS 3DEP / Illinois Height Modernization Program covers Cook County at 1--8 pts/m$^2$; Microsoft GlobalMLBuildingFootprints + USGS 2023 building-heights release explicitly include Chicago. A handheld rangefinder is useful only for spot-checking the 5--10 buildings nearest the field.
  \item \textbf{``OpenFOAM or COMSOL''} --- pick OpenFOAM. The urban-CFD literature (Blocken, Tominaga, COST 732 community) is overwhelmingly OpenFOAM-based; \texttt{snappyHexMesh} handles dirty STL of campus buildings well; you'll have peer-reviewed cases to validate against.
\end{itemize}

\subsection{The biggest unstated gap}

You do not have a \emph{bridging quantity} that turns a CFD flow field into a throw-quality score. The paper needs one. \S\ref{sec:throwability} proposes a ``throwability metric'' $Q(\mathbf{x}, \theta_{\text{wind}})$ combining mean wind magnitude, turbulence intensity at disc altitude, and predicted lateral trajectory drift for the canonical throws.

\section{Project plan}
\label{sec:plan}

\subsection{Work breakdown ($\approx$14 weeks, single researcher)}

\begin{center}
\footnotesize
\begin{tabular}{@{}llp{0.55\linewidth}@{}}
\toprule
\textbf{Phase} & \textbf{Weeks} & \textbf{Output} \\
\midrule
P0 --- Setup & 1 & OpenFOAM 11 or v2406 installed; Python disc-trajectory simulator stub; QGIS workflow for LiDAR $\to$ STL. \\
P1 --- Analytical disc model & 2--3 & 6-DOF integrator (Python, scipy \texttt{solve\_ivp}). Validated against Hummel (2003) Fig.\ 5-1 throw reconstructions. \\
P2 --- Geometry build & 4 & STL of $\sim$250 m radius around Sean Earl Field, $\sim$5--10 named buildings + tree-canopy proxy. \\
P3 --- RANS CFD sweep & 5--8 & 13 steady-state runs (15$^\circ$ increments). \texttt{simpleFoam} + k-$\omega$ SST + Richards--Hoxey ABL inlet. Output: $U_{\text{mean}}, k$ sampled to 1, 2, 3 m AGL. \\
P4 --- Coupled trajectory & 9--11 & Trilinear-interpolated CFD into 6-DOF disc model. Monte Carlo over launch, heading, spin, release $\alpha$ ($10^4$ throws per direction). \\
P5 --- Synthesis \& manuscript & 12--14 & Throwability map $Q(\mathbf{x}, \theta_{\text{wind}})$; paper revision; figures. \\
\bottomrule
\end{tabular}
\end{center}

\subsection{Compute and storage}

Workstation target: 32 physical cores, 128 GB RAM, $\sim$2 TB scratch SSD. On 16-core / 64 GB expect $\sim$2$\times$ wall-clock and one wind direction per overnight job. Disk budget: $\sim$50 GB per RANS case ($\sim$700 GB for the sweep). LES would multiply by 20--50$\times$.

\subsection{Inputs (all free/accessible)}

\begin{itemize}[leftmargin=*,itemsep=2pt]
  \item \emph{Building footprints/heights:} Cook County GIS, Microsoft GlobalMLBuildingFootprints (ODbL), USGS Overture Maps 2023, OpenStreetMap.
  \item \emph{LiDAR DEM:} USGS 3DEP / Illinois Height Modernization Program (1--8 pts/m$^2$).
  \item \emph{Wind reference:} NWS O'Hare 1991--2020 normals; NOAA GLERL Wilson Avenue Crib (shoreline calibration); Northerly Island ASOS.
  \item \emph{Disc aerodynamic coefficients:} Hummel (2003) Tables 3-1/3-2; Crowther \& Potts (2007).
\end{itemize}

\subsection{Risks \& mitigations}

\begin{center}
\footnotesize
\begin{tabular}{@{}p{0.46\linewidth}p{0.46\linewidth}@{}}
\toprule
\textbf{Risk} & \textbf{Mitigation} \\
\midrule
k-$\omega$ SST under-predicts rooftop separation at sharp roof corners (well-known). & Validate one direction against a LES spot check; follow Blocken (2018) on RANS limits. Don't claim peak/fluctuating quantities --- report means + TKE only. \\
Disc-coefficient uncertainty at off-design $\alpha$ ($>20^\circ$). & Restrict throw library to $\alpha \in [-5^\circ, +20^\circ]$; flag hammer (high-$\alpha$) results as exploratory. \\
Horizontal inhomogeneity of ABL inlet (Hargreaves \& Wright 2007). & Use Yang et al.\ (2009) corrected wall functions; verify inlet profile preserved upstream of buildings. \\
Sean Earl is not a competition-grade Ultimate field. & Reframe as a methods paper using Sean Earl as exemplar; methodology generalizes. \\
Reviewers ask for measurement validation. & Use Tominaga (2008) AIJ-B cube benchmark for the CFD; cite Hummel's marker-tracked flights for the disc model. Both are public. \\
\bottomrule
\end{tabular}
\end{center}

\section{Abstract (draft)}
\label{sec:abstract}

Outdoor disc sports are dominated by wind, yet the published aerodynamic and trajectory literature for flying discs (Hummel 2003; Potts \& Crowther 2002; Lukes et al.\ 2022) uses uniform-freestream assumptions that ignore site-specific flow structure. Conversely, the urban-CFD literature (Tominaga et al.\ 2008; Blocken 2014, 2018) characterizes site-specific wind fields but stops at pedestrian-level mean quantities. Here we couple a validated 6-degree-of-freedom (6-DOF) disc-flight model to steady-RANS atmospheric CFD of a lakefront university campus (Loyola University Chicago, Sean Earl Field; $\sim$200 m from the Lake Michigan shore). We sweep 13 inflow directions in 15$^\circ$ increments using OpenFOAM with k-$\omega$ SST turbulence and a Richards--Hoxey atmospheric-boundary-layer inlet, validated against COST Action 732 benchmarks. Disc trajectories for canonical throws (backhand, forehand, hammer; 175~g Ultrastar) are integrated through the resulting time-averaged + turbulence-intensity fields. We report a spatial ``throwability'' map $Q(\mathbf{x}, \theta_{\text{wind}})$ and identify the lakefront internal boundary layer --- formed by the sharp roughness transition at the shoreline --- as the dominant local feature controlling disc-altitude flow. The methodology generalizes to any outdoor sport played on a roughness-discontinuity site.

\section{Introduction (draft)}

The Frisbee is one of the most-studied non-airfoil aerodynamic objects in sports physics. Stilley \& Carstens (1972) provided the first parametric model; Yasada (1999), Hummel (2003), and Potts \& Crowther (2002) supplied successively more complete wind-tunnel measurements of the lift, drag, and pitching-moment coefficients, and Hummel's UC Davis thesis remains the most-cited 6-DOF flight reconstruction. Subsequent CFD validation by Kamaruddin (2011) and Lukes et al.\ (2022, 2023) using RANS k-$\omega$ SST in OpenFOAM has converged the modeling stack to within a few percent of wind-tunnel $C_L$, $C_D$, and $C_M$ at moderate angles of attack.

Yet every published disc-trajectory study assumes a spatially uniform, steady freestream. As Lukes et al.\ (2022) explicitly note, ``the biggest uncertainty in trajectory experiments is the wind environment.'' This uncertainty is not noise to be averaged away: outdoor disc sports --- Ultimate, disc golf --- are played in geographically structured turbulent flow. A throw whose intended landing zone lies in a sheltered eddy behind a building experiences a categorically different flow than the same throw 20 m to the side.

Independently, the urban-CFD community has developed mature best-practice guidance for site-specific pedestrian-level wind: the COST Action 732 guidebook (Franke et al.\ 2007), the AIJ guidelines (Tominaga et al.\ 2008), and the 50-year retrospective of Blocken (2014). Standard practice uses steady RANS with k-$\omega$ SST or realizable k-$\varepsilon$ on \texttt{snappyHexMesh}-generated meshes of $\sim 10^7$ cells, with Richards--Hoxey atmospheric inlets. This work has shaped pedestrian-comfort regulation in dozens of cities, but to our knowledge has never been applied to outdoor sports trajectory prediction.

The present work fills that gap. We choose Sean Earl Field at Loyola University Chicago's Lake Shore Campus as a study site because it offers an unusually clean physical setup: the field sits within $\sim$250 m of Lake Michigan, with no major intervening built obstacle to the east, behind it a typical low-rise urban-campus block to the west. Easterly winds therefore arrive after $\sim$80 km of unbroken over-water fetch ($z_0 \approx 0.001$ m), pass through a sharp roughness transition at the shoreline, and develop a shallow internal boundary layer over the field. Westerly winds arrive through the campus block. The site cleanly separates two regimes (smooth-fetch and urban-fetch) within a 180$^\circ$ heading sweep.

\textbf{Our contributions:}
\begin{enumerate}[leftmargin=*,itemsep=2pt]
  \item The first published two-way coupling of an environmental-CFD wind field to a validated 6-DOF disc-flight model.
  \item A ``throwability metric'' $Q(\mathbf{x}, \theta_{\text{wind}})$ that compresses the spatial CFD field plus trajectory uncertainty into an actionable single scalar per (launch position, wind heading).
  \item A case study identifying the lakefront internal boundary layer as the dominant local control on Ultimate-disc flight at the chosen site.
\end{enumerate}

\section{Analytical disc aerodynamics (summary)}

Full derivation in the companion document. The disc carries a 6-DOF state $(\mathbf{x}, \mathbf{v}, \psi, \theta, \phi_s, p, q, \Omega)$ governed by Newton's equation (\ref{eq:trans-summary}) plus Euler's equation for a symmetric rotor (\ref{eq:rot-summary}). In the high-spin limit relevant to all real throws, the transverse angular rates are slaved algebraically to the applied aerodynamic moments,
\begin{equation}
p = -\frac{M}{I_p \Omega}, \qquad q = \frac{R}{I_p \Omega},
\end{equation}
recovering the textbook gyroscopic-precession picture. The pitching moment vanishes at the trim angle
\begin{equation}
\alpha_{\text{trim}} = -\frac{C_{M0}}{C_{M\alpha}} = 8.8^\circ \text{ (Potts \& Crowther coeffs)},
\label{eq:alphatrim}
\end{equation}
which is the asymptotic glide angle of attack. The full translational and rotational equations entering the numerical simulator are
\begin{align}
m\,\dot{\mathbf{v}} &= -mg\,\hat{\mathbf{n}}_3 + \tfrac{1}{2}\rho V^2 A\big[\,C_L(\alpha)\hat{\mathbf{L}} - C_D(\alpha)\hat{\mathbf{v}}_{\text{rel}}\big], \label{eq:trans-summary}\\
I_p \dot{p} + I_p q \Omega &= R,\qquad I_p \dot{q} - I_p p \Omega = M,\qquad I_s \dot{\Omega} = N. \label{eq:rot-summary}
\end{align}

\paragraph{Coefficients and their functional form.} We use the tabulated
Potts \& Crowther (2002) $C_L$, $C_D$, $C_M$ curves and the
$C_{Rr}(\alpha, \text{AdvR})$ surface as transcribed in Hummel (2003,
Appendix A), evaluated exactly as in Hummel's reference implementation:
$C_L = C_{L0} + C_{L\alpha}\alpha$ and $C_M = C_{M0} + C_{M\alpha}\alpha$ are
functions of $\alpha$ itself, while the zero-lift angle
$\alpha_0 = -C_{L0}/C_{L\alpha}$ enters \emph{only} the drag parabola
$C_D = C_{D0} + C_{D\alpha}(\alpha - \alpha_0)^2$. The rate-damping
coefficients $C_{Mq}$, $C_{Rp}$ multiply raw angular rates in rad\,s$^{-1}$ and
$C_{Rr}$ carries an explicit $\sqrt{d/g}$ time scale; they are not
non-dimensionalised by $d/2V$. Mass and inertia are Hummel's measured Table 3-1
values ($m = 0.175$ kg, $I_s = 2.35\times10^{-3}$, $I_p = 1.22\times10^{-3}$
kg\,m$^2$), not the thin-uniform-disc idealisation, which under-states both by
about one third and so over-states precession by roughly 50\%.

\paragraph{Validation.} The integrator is checked against two independently
published, fully specified results in the same source that supplied the
coefficients (Table~\ref{tab:p1valid}). Hummel's \S2.4.3 ``common throw''
(release $V_0=14$ m\,s$^{-1}$, $\Omega=50$ rad\,s$^{-1}$, pitch $11^\circ$,
flight-path angle $6^\circ$, $\alpha=5^\circ$, $z_0=1$ m) is reproduced within
8\% in range, 4\% in apex height and 1\% in flight time, and the model
reproduces the characteristic turn-right-then-fade-left signature with the
reversal at $t=3.34$ s against a published $\sim$3 s. The \S2.4.2 analytic trim
glide is recovered to three significant figures.

\begin{table}[H]
\centering
\small
\begin{tabular}{llll}
\toprule
Quantity & Hummel (2003) & This model & Error \\
\midrule
\multicolumn{4}{l}{\emph{\S2.4.3 / Fig. 2-9 ``common throw''}}\\
Range                       & $\sim$25 m   & 26.99 m  & $+8\%$ \\
Apex height                 & 2.55 m       & 2.45 m   & $-4\%$ \\
Flight time                 & 3.5 s        & 3.48 s   & $-1\%$ \\
Terminal speed              & $\sim V_0/3$ & 4.83 m\,s$^{-1}$ & $+4\%$ \\
Lateral reversal            & $\sim$3 s    & 3.34 s   & --- \\
\multicolumn{4}{l}{\emph{\S2.4.2 steady glide (analytic)}}\\
Trim $\alpha$ ($C_M=0$)     & $9^\circ$    & $9.00^\circ$   & $<0.1\%$ \\
Glide angle $\gamma$        & $-19.3^\circ$& $-19.32^\circ$ & $<0.2\%$ \\
$C_L$ at trim               & 0.67         & 0.672    & $+0.3\%$ \\
Trim speed                  & 9.1 m\,s$^{-1}$ & 8.28 m\,s$^{-1}$ & $-9\%$ \\
\bottomrule
\end{tabular}
\caption{Quantitative validation of the 6-DOF integrator against published
reference flights. Only the tabulated coefficient set reproduces these; the
linear fits miss range, apex and flight time by 30--57\% and never fade back
left, because a realistic throw reaches $\alpha \approx 33^\circ$ in its final
second, far outside the $\pm15^\circ$ band in which the linear fits were made.}
\label{tab:p1valid}
\end{table}

\paragraph{Known limitations of the disc model.} The model is axisymmetric and
carries no rim side-force coefficient $C_S(\alpha)$, hence no yaw stiffness.
Launch sideslip is therefore shed only slowly, and predicted crosswind drift is
about a quarter of the naive full-advection estimate $U_c T$. We verified that
this is a physics-model gap rather than a numerical artefact: a uniform wind is
exactly Galilean-equivalent to a shifted release velocity in this
implementation, to $10^{-8}$ m. \textbf{Crosswind drift from this model should
be read as a lower bound.}

\section{Site description: Sean Earl Field}

Sean Earl Field is a small multi-use turf on Loyola's Lake Shore Campus, with
Mertz Residence Hall to the west and the Norville Center for Intercollegiate
Athletics to the south. Loyola explicitly lists Ultimate among the field's
sanctioned uses.

\paragraph{Geometry as built, and four corrections to the initial description.}
The site model is built from OpenStreetMap footprints with
\texttt{building:levels} heights, with the local Cartesian origin resolved at
run time from the OSM polygon of the field itself
($42.001267^\circ$ N, $87.658214^\circ$ W; $x$ East, $y$ North). Building the
model surfaced four discrepancies against the description we began with, each
of which changes the setup:

\begin{enumerate}[leftmargin=*,itemsep=2pt]
  \item \textbf{Origin.} The initially assumed centre
    ($42.0006^\circ$ N, $87.6580^\circ$ W) is $\sim$75 m south of the field
    centroid; using it would have offset the domain relative to the launch grid.
  \item \textbf{Mertz Hall is 19 storeys, $\approx$61 m, not $\approx$22 m}
    (OSM \texttt{building:levels}\,$=$\,19, corroborated by Loyola's own
    residence-life page). This matters directly: domain sizing under COST
    Action 732 scales with the tallest building.
  \item \textbf{The field's long axis runs E--W, not N--S.} The OSM polygon is
    78 m (E--W) $\times$ 25 m (N--S). \emph{\% FLAG FOR AUTHOR: this is also
    smaller than a regulation Ultimate pitch (100 $\times$ 37 m), so the OSM
    outline is likely a partial trace of the playing surface. The launch grid
    below is sized to the OSM polygon; a corrected outline would extend it.}
  \item \textbf{Cudahy Library and the Klarchek Information Commons lie ESE of
    the field} (129 m and 185 m respectively), on the lakefront side, not to
    the west as first described.
\end{enumerate}

\paragraph{The dominant near-field obstruction.} The single most consequential
feature is one the initial site description omits entirely. \textbf{Campus
Towers}, a 14-storey ($44.8$ m) apartment block, occupies
$x \in [55.8, 79.3]$ m, $y \in [-15.9, 21.5]$ m: its footprint begins
\textbf{9 m east of the field's eastern edge} and spans the full width of the
field. Any easterly flow reaching the playing surface has therefore passed
over a 45 m bluff body immediately upwind. \S\ref{sec:results} shows that this,
not the lakefront roughness transition, controls disc-altitude flow in the
easterly regime.

\paragraph{Terrain.} USGS 3DEP point queries over a 600 m box give
179.1--181.4 m MSL across the field and the near-field buildings ($\pm 1$ m),
against building heights of 12--61 m, so the ground is modelled as flat. The
one real topographic feature is a $\sim$4 m step down to the lake east of
$x \approx +150$ m (3DEP reads 175.2--175.6 m there); it is represented not as
topography but where it acts, as the inlet roughness transition.

\paragraph{Climatology.} Chicago lies in the mid-latitude westerly belt; the annual prevailing direction at ORD is WSW--SW (Illinois State Water Survey wind roses). The 1991--2020 NWS annual mean speed at O'Hare is $\sim$9--10 mph. However, ORD is 27 km inland; lakefront stations (NOAA GLERL Wilson Avenue Crib, Northerly Island ASOS) record higher mean speed and gustier easterlies, especially during the warm-season lake-breeze regime.

\paragraph{Why this site.} Easterly inflow at Sean Earl Field passes over $\sim$80 km of unbroken Lake Michigan fetch (to St.\ Joseph / Benton Harbor MI), giving a near-neutral, very-low-$z_0$ marine boundary layer. The shoreline imposes a sharp $z_0$ transition ($\sim 10^{-3}$ to $\sim$0.3 m). The downwind internal boundary layer thickness $\delta_{\text{IBL}}$ grows as $\delta_{\text{IBL}} \approx 0.10\,x^{0.8}$ (Garratt 1990), reaching $\sim$5 m by the eastern sideline and $\sim$7--10 m at midfield --- fully encompassing the 1--3 m disc-altitude band but in a sheared, non-equilibrium state. \emph{This is the dominant local physical feature.}

\section{CFD methodology}

\paragraph{Geometry.} Building footprints come from OpenStreetMap via the
Overpass API, extruded to \texttt{building:levels} heights (171 of 303 fetched
buildings, 56\%, carry a levels tag; the remainder use documented type
defaults, and each building records which rule was used). The modelled set is
clipped to a \emph{circle} of radius 300 m about the field centroid: this makes
the building set invariant under the rotation used to realise each wind
direction (below), so every direction sees exactly the same geometry. Within
that circle the tallest building is Mertz Hall, $H = 60.8$ m. Ground is flat
(\S6) and no tree-canopy porous zone is included; both omissions are recorded
in \S\ref{sec:limits}. The resulting surface is 116 closed prisms, 6\,324
triangles, with no open edges.

\paragraph{Domain.} $1512 \times 1208 \times 304$ m, i.e.\ $5H$ upstream,
$5H$ lateral and $5H$ above the tallest building, and $10H$ downstream.
The downstream extent is \textbf{relaxed} from the COST Action 732 /
Tominaga et al.\ (2008) guideline of $15H$ to fit the available compute; the
relaxation is recorded in every case's metadata rather than left implicit.
Blockage is $0.9\%$, within the $<3\%$ guideline. (For reference, the
$500\times500\times80$ m domain in the original plan would have placed the
domain top \emph{below} the tallest building.)

\paragraph{Mesh.} \texttt{snappyHexMesh}, hex-dominant, on a graded background
of 16 m horizontal cells with a 4 m first vertical cell expanding geometrically
to the domain top. Two refinement regions: a near-field box at level 1 and a
disc-altitude box over the playing surface at level 3, giving
\textbf{2.0 m horizontal $\times$ 0.50 m vertical} cells through the
1--6 m AGL band the results are sampled in; building surfaces are refined to
levels 1--2. Production meshes are $4.5\times10^{5}$ cells. Surface layers are
\emph{not} added: with rough-wall ABL functions there is no viscous sublayer to
resolve, and layer insertion on 116 real footprints is the least robust meshing
step. Mesh quality: maximum non-orthogonality 65$^\circ$ (mean 7.9$^\circ$),
maximum skewness 7.9.

\paragraph{Solver.} OpenFOAM \textbf{v2606} (ESI) \texttt{simpleFoam}, steady
incompressible, k-$\omega$ SST, second-order schemes with limited
($\psi = 0.33$) Laplacian and surface-normal-gradient corrections for the
non-orthogonal cells snapping to real footprints produces. Richards \& Hoxey
(1993) ABL inlet (\texttt{atmBoundaryLayerInletVelocity}, \texttt{...InletK},
\texttt{...InletOmega}) with the atmospheric rough-wall functions
(\texttt{atmNutkWallFunction}, \texttt{atmOmegaWallFunction}) on the ground.
Reference $u_{\text{ref}} = 5$ m\,s$^{-1}$ at $z_{\text{ref}} = 10$ m.

\paragraph{Roughness, and why the ground and the inlet differ.} Inlet
roughness follows the plan: $z_0 = 0.001$ m (open water) for
$\theta_{\text{wind}} \in [60^\circ, 120^\circ]$, $z_0 = 0.5$ m (urban)
otherwise, linearly blended over $30^\circ$ transition bands. The
\emph{ground patch}, however, uses $z_0 = 0.03$ m (grass and pavement), not the
urban 0.5 m. The 0.5 m figure is an \emph{equivalent} roughness that stands in
for buildings; applying it to the ground as well would count the resolved
buildings twice. It is also numerically incompatible with the mesh: the ABL
wall function requires the first cell centre to sit well above $z_0$, and a
0.5 m roughness cannot coexist with the sub-metre near-ground cells needed to
resolve disc altitude.

\paragraph{Realising the wind directions.} Rather than moving the inlet between
patches --- which only works for the four axis-aligned headings and gives each a
different mesh --- the \emph{site geometry is rotated} so the wind always blows
along $+x$. Every direction then shares one background mesh, one boundary-
condition set and one inlet patch, so differences between directions are
physical rather than mesh artefacts. This matters because the paper's claims
are comparisons \emph{between} directions.

\paragraph{Cases actually run.} \textbf{Three inflow directions, not the
thirteen originally planned}, chosen to span the three regimes the study
hypothesises: $\theta_{\text{wind}} = 90^\circ$ (easterly, over-lake),
$210^\circ$ (SSW, the prevailing sector), and $270^\circ$ (westerly, over the
campus and Mertz Hall). Note that $270^\circ$ extends past the
$0^\circ$--$180^\circ$ range of the original plan; the geometry-rotation scheme
makes any heading equally cheap. Steady RANS, 2\,500 iterations, at a single
reference wind speed. \S\ref{sec:limits} states plainly what this subset does
and does not license.

\paragraph{Validation.} Two of the three planned checks were carried out and
are reported in \S\ref{sec:cfdvalid}: ABL horizontal homogeneity over an empty
fetch (Hargreaves \& Wright 2007) and grid independence at disc altitude
against the 10\% criterion. \textbf{The Tominaga (2008) AIJ-B cube benchmark
was not run}, and no anemometer comparison exists; neither is claimed.

\section{Coupled trajectory simulation}

For each $\theta_{\text{wind}}$, the CFD output is sampled onto a 1 m horizontal
grid at $z = 1$--6 m AGL across the field and trilinearly interpolated. The
6-DOF integrator evaluates $\mathbf{U}(\mathbf{x})$ at the disc's instantaneous
position; the steady-RANS field provides $U_{\text{mean}}$, and the local
turbulence intensity $I = \sqrt{2k/3}/U_{\text{mean}}$ scales a stochastic gust
$\mathbf{u}'(t)$.

\paragraph{Gust model, and why it is not white noise.} The gust field is
generated by a Shinozuka--Deodatis spectral representation: a sum of 32 cosines
with log-spaced frequencies and independent random phases, with amplitudes set
by an exponential-correlation spectrum whose integral time scale is
$T_L = k/\varepsilon$ from the RANS, and normalised to unit variance so that
$\mathbf{u}'(\mathbf{x},t) = \sigma(\mathbf{x})\,\bm{\xi}(t)$ with
$\sigma = \sqrt{2k/3}$.

The obvious cheaper choice --- drawing an independent Gaussian fluctuation each
time the wind is queried --- is not merely inaccurate but unusable, and the
reason is worth stating because it is easy to reintroduce. An adaptive ODE
integrator estimates local truncation error by comparing right-hand-side
evaluations at different stage points within a step. If the right-hand side
returns fresh noise on every call, that comparison measures the noise rather
than the discretisation error, so the step controller rejects and halves
indefinitely. In our tests a single throw took over 100 s against $\sim$0.1 s
without gusts, and the resulting landing scatter is a step-size artefact rather
than a solution of any well-posed equation. The spectral construction is
analytic in $t$, so the integrator steps freely (193 steps with gusts against
222 without) and the coupled system is well posed.

The realised autocorrelation matches the target $\exp(-\Delta t/T_L)$ to within
0.06 over lags out to 2 s. The construction is the same family as the Mann
(1998) sampler; upgrading amounts to substituting the spectrum.

For each $(\theta_{\text{wind}}, \text{throw type})$, $10^4$ Monte Carlo throws are launched from a 5$\times$5 m grid of launch points spanning the field with realistic release-parameter scatter:
\begin{itemize}[leftmargin=*,itemsep=1pt]
  \item Speed $V_0 \sim \mathcal{N}(15, 1.5)$ m/s for backhand and forehand; $\mathcal{N}(11, 1.5)$ for hammer.
  \item Spin $\Omega_0 \sim \mathcal{N}(100, 10)$ rad/s.
  \item Release pitch $\alpha_0 \sim \mathcal{N}(+8^\circ, 4^\circ)$ backhand; $\mathcal{N}(+6^\circ, 4^\circ)$ forehand; $\mathcal{N}(+40^\circ, 8^\circ)$ hammer.
  \item Heading bearing $\theta_{\text{throw}}$ varied over 0--360$^\circ$ in 30$^\circ$ bins.
\end{itemize}
Landing-point distributions are characterized by 90\%-containment ellipse parameters and a ``delivered-to-target'' success rate (landing within 3 m of intended target at canonical 25 m throws).

\section{Throwability metric (the novelty)}
\label{sec:throwability}

For each launch position $\mathbf{x}$ and each wind heading $\theta_{\text{wind}}$:
\begin{equation}
Q(\mathbf{x}, \theta_{\text{wind}}) =
\underbrace{\Big(1 - \frac{|U_{2\text{m}}(\mathbf{x})|}{U_{\text{ref}}}\Big)}_{\text{wind-magnitude penalty}}
\cdot
\underbrace{\big(1 - I_{2\text{m}}(\mathbf{x})\big)}_{\text{turbulence penalty}}
\cdot
\underbrace{\overline{S}_{\text{throws}}(\mathbf{x}, \theta_{\text{wind}})}_{\text{simulated success rate}}
\label{eq:Q}
\end{equation}
with each factor clipped to $[0, 1]$. $Q$ ranges from 0 (unplayable) to 1
(calm-air-equivalent). Maps of $Q$ over the field surface, parameterised by
$\theta_{\text{wind}}$, form the principal result.

\paragraph{Two points on the definition.} First, the wind-magnitude penalty uses
the \emph{local} $|U_{2\text{m}}(\mathbf{x})|$. An earlier form of this metric
used a field-averaged $\langle |U_{2\text{m}}|\rangle_{\text{field}}$, which is
constant in $\mathbf{x}$ and therefore cannot contribute any spatial structure
at all --- it can only rescale a map, not shape one. Since the spatial gradient
of $Q$ is the result being reported, the local value is the right choice.

Second, $\overline{S}_{\text{throws}}$ is the \emph{measured} fraction of Monte
Carlo throws landing within 3 m of a 25 m target, computed from the simulated
landing points. It is tempting to estimate it from the containment ellipse as
$\min\!\big(1, (r_{\text{target}}/\bar a)^2\big)$; that quantity is not a
probability, because it ignores the offset between the landing centroid and the
target and so scores a tightly grouped but badly aimed throw as a success.

\paragraph{A caveat on the turbulence factor.} The form $(1 - I)$ goes negative
wherever the RMS fluctuation exceeds the local mean speed. That is not an exotic
condition: it is the normal state of a recirculating building wake, which is
exactly the sheltered region a throwability map is most interesting in.
Clipping at zero is the only defensible choice --- a calm sheltered spot must
not score worse than a gale --- but it means $Q$ collapses to identically zero
across the whole saturated region and carries no spatial information there. We
report the saturated area fraction for every heading in \S\ref{sec:results}
rather than letting it disappear into the clip, and we recommend a bounded
penalty such as $1/(1+I)$ for sheltered sites in future work.

\section{CFD validation}
\label{sec:cfdvalid}

\paragraph{(a) ABL horizontal homogeneity.} An equilibrium boundary layer
imposed at the inlet should advect across an empty domain unchanged; any drift
is spurious and contaminates everything downstream (Hargreaves \& Wright 2007).
We ran a buildings-free case with the ground roughness set equal to the inlet
roughness --- necessary, because with the two deliberately different (as in the
production cases, where the mismatch \emph{is} the shoreline roughness
transition) the profile is supposed to evolve and the test would measure the
intended physics rather than numerical error.

Profiles were compared at five fetches from 50 m to 1400 m, evaluated above the
first wall-function cell centre. The worst streamwise drift in mean velocity is
\ABLdrift\% (\ABLverdict), and the profiles depart from the analytic
Richards--Hoxey curve by 1.3\%; turbulent kinetic energy holds to within about
7\% of $u_*^2/\sqrt{C_\mu}$. Sampling below the first cell centre is not
meaningful with wall functions --- it interpolates toward the no-slip wall
value --- and is excluded.

One caveat on the reach of this test. The empty-fetch case has no buildings, so
\texttt{snappyHexMesh} is skipped and the mesh is the unrefined background,
whose first cell is 4 m tall. The test therefore establishes that the
inlet/wall-function/turbulence-model stack is mutually consistent \emph{aloft};
it does not by itself certify the profile at 1--6 m AGL. In the production cases
that band is resolved by roughly twelve 0.5 m cells inside the refinement box,
which is the relevant configuration, but a dedicated refined empty-fetch case
would be the cleaner test and is left as follow-on work.

\paragraph{(b) Grid independence.} Two meshes differing \emph{only} in the
disc-altitude refinement level (2.0 $\times$ 0.50 m against 4.0 $\times$ 1.0 m,
a factor of two in the region we report) were compared at five stations over
the playing surface. The worst disc-band mean-velocity difference is
\GridDiff\% of the local speed, or \GridDiffUref\% of the reference speed
$u_{\text{ref}} = \Uref$ m\,s$^{-1}$ (\GridVerdict).

Both measures are quoted deliberately. The field turns out to be heavily
sheltered, with disc-altitude speeds of a few tenths of a metre per second in
places, so a purely relative criterion is punishing there: a 0.3 m\,s$^{-1}$
absolute difference is a few per cent of $u_{\text{ref}}$ but can be a large
fraction of the local speed. Neither number should be read as a clean pass, and
\S\ref{sec:limits} states the consequence.

\paragraph{(c) Convergence, and why the residual criterion is the wrong test.}
Steady SIMPLE does not reach a fixed point on this geometry. The playing surface
sits in the separated near-wake of a 45 m building; the physical flow there is
unsteady, and the solver settles into a limit cycle in which the pressure
residual plateaus near $5\times10^{-3}$ against the $10^{-4}$ target that the
project plan specified. Reporting that target as met would be false; reporting
the run as simply ``unconverged'' would be equally misleading, because the
residual has stopped falling for a physical reason rather than a numerical one.

We therefore assess convergence on the quantity we actually report. Velocity
probes at disc altitude over the field oscillate by \ConvFine\% peak-to-peak
over the last fifth of the run on the production mesh. Because a single iterate
is an arbitrary phase of that cycle, all reported fields are \textbf{time
averages} of $\mathbf{U}$, $k$ and $\omega$ over a further 800 iterations after
the cycle is established, not final iterates.

\section{Expected results / placeholder findings}

To be populated after the simulation sweep. Hypothesized:
\begin{itemize}[leftmargin=*,itemsep=2pt]
  \item Easterly cases ($90^\circ \pm 30^\circ$) show a sharp $Q$ gradient across the field's eastern half corresponding to the developing IBL; backhand throws from the windward sideline are most consistent; hammer throws systematically degraded.
  \item Westerly cases ($270^\circ \pm 30^\circ$) show patchier $Q$ with sheltered ``throw lanes'' downstream of Norville Center; forehand throws benefit most.
  \item Best-condition wind direction for canonical 25 m throws is hypothesized to be a light SSW ($\sim$210$^\circ$, 3--5 m/s reference), aligning with prevailing climatology --- a satisfying convergence of meteorology and playability.
\end{itemize}

\section{Discussion (to draft after results)}

Generalizability to other outdoor sports; steady-RANS limitations (peak gusts and spatially correlated gust structure not captured); LES of one or two critical directions is the obvious follow-on; coaching implications.

\section{Scope and limitations of what was actually computed}
\label{sec:limits}

This section states the delivered scope plainly, so that no reader has to infer
it from the methods.

\paragraph{CFD directions.} \textbf{Three} inflow directions were computed:
$\theta_{\text{wind}} = 90^\circ$, $210^\circ$ and $270^\circ$. The
thirteen-direction, $15^\circ$-increment sweep described in the project plan
(\S\ref{sec:plan}) \textbf{was not run}. The three directions were chosen to
span the three regimes the study hypothesises --- over-lake easterly,
prevailing SSW, and over-campus westerly --- and they are sufficient to
\emph{contrast} those regimes, but they cannot resolve the direction-dependence
of $Q$ between them. Any statement in this paper about a ``best'' or ``worst''
heading means best or worst \emph{of the three computed}, and is written that
way throughout.

\paragraph{Mesh resolution.} Production meshes are $4.5 \times 10^{5}$ cells,
not the $5$--$8 \times 10^{6}$ of the plan. The refinement is concentrated where
the results are read: $2.0 \times 0.50$ m cells through the 1--6 m AGL band over
the playing surface. Grid independence is demonstrated at that resolution
(\S\ref{sec:cfdvalid}) against a mesh a factor of two coarser in the same
region, which bounds discretisation error where it matters; it does not
establish that far-field wake structure is grid-converged.

\paragraph{Domain.} Upstream, lateral and top extents meet the $5H$ COST 732
guideline; the downstream extent is $10H$ against a $15H$ guideline. Blockage is
$0.9\%$.

\paragraph{Wind speed.} A single reference speed,
$u_{\text{ref}} = 5$ m\,s$^{-1}$ at 10 m. Because the flow is fully turbulent
and the geometry is bluff, the normalised flow field is close to
Reynolds-independent, so the $Q$ maps rescale approximately with
$u_{\text{ref}}$ --- but the trajectory model is \emph{not} linear in wind
speed, so throwability at 10 m\,s$^{-1}$ cannot be read off these results.

\paragraph{Validations run and not run.} ABL horizontal homogeneity and grid
independence were performed and are reported. The \textbf{Tominaga (2008) AIJ-B
cube benchmark was not run}. \textbf{No field measurements were made}, so
nothing here is validated against the real site; the anemometer campaign
remains future work.

\paragraph{Physical modelling.} Steady RANS gives mean-condition playability
only: no peak gusts, and no spatially correlated gust structure. The gust model
is temporally correlated but spatially uniform along the trajectory at any
instant, isotropic (RANS $k$ carries no anisotropy), and Gaussian, so extreme
gusts are under-represented. The disc model has no rim side-force term, so
predicted crosswind drift is a lower bound. Ground is flat and no tree canopy is
included. The hammer throw is exploratory: it spends most of its flight outside
the angle-of-attack range over which the tabulated coefficients exist.

\section{Author decisions}
\label{sec:decisions}

The open questions carried by the draft are resolved below with defaults that
the present results support. Each remains the author's to overrule.

\begin{enumerate}[leftmargin=*,itemsep=3pt]
  \item \textbf{Site.} Sean Earl Field, Loyola Lake Shore Campus, is assumed
    and used throughout, with the geometry resolved from OSM at
    $42.001267^\circ$ N, $87.658214^\circ$ W. \emph{\% FLAG FOR AUTHOR: confirm
    the field outline, which OSM traces as 78 $\times$ 25 m --- smaller than a
    regulation Ultimate pitch.}

  \item \textbf{LES.} Out of scope here and recorded as the primary follow-on.
    The steady-RANS result gives mean-condition playability; it cannot give
    peak gusts or spatially correlated gust structure, both of which matter for
    a thrower. An LES of the single easterly case would test precisely the
    near-wake claim in \S\ref{sec:results} that RANS is least reliable about.
    \emph{\% FLAG FOR AUTHOR: $\sim$2 weeks of compute.}

  \item \textbf{Anemometer validation.} Stated as a future field campaign, not
    claimed. No field measurements were available for this work, and none are
    reported. \emph{\% FLAG FOR AUTHOR: if the 2026 summer lake-breeze campaign
    goes ahead, the three field-corner stations proposed in the draft are the
    right comparison points; the CFD prediction at those points is already in
    \texttt{flow\_fields.npz}.}

  \item \textbf{Hammer throw.} Retained but reported as \textbf{exploratory}.
    It is the throw most affected by wind and the one whose coefficients are
    least well supported: a hammer spends most of its flight well outside the
    $-10^\circ$ to $+30^\circ$ band over which the tabulated Potts \& Crowther
    coefficients exist, so its trajectories rely on edge-clamped values. Its
    numbers are reported separately from backhand and forehand and are not used
    in any headline claim.

  \item \textbf{Authorship and affiliation.}
    \emph{\% FLAG FOR AUTHOR: author name, department, affiliation, co-authors
    and funding acknowledgement all require the author's input and are left as
    placeholders on the title page.}

  \item \textbf{Venue.} \emph{\% FLAG FOR AUTHOR:} on the strength of the
    present result --- a single-site coupled study with a scoped CFD sweep ---
    \emph{Sports Engineering} remains the better fit than splitting into
    separate trajectory and wind-engineering papers. The disc-model validation
    (Table~\ref{tab:p1valid}) is solid enough to stand alone, but the CFD side
    is currently three directions at one reference wind speed, which is thin
    for \emph{JWEIA} on its own. Splitting becomes attractive after the full
    direction sweep and the LES spot-check.
\end{enumerate}

\section*{Key references (verified)}

\begin{description}[font=\normalfont\itshape,leftmargin=1em]
  \item Blocken, B.\ (2014). 50 years of computational wind engineering: Past, present and future. \emph{JWEIA} \textbf{129}, 69--102.
  \item Blocken, B.\ (2018). LES over RANS in building simulation\ldots\ \emph{Building Simulation} \textbf{11}, 821--870.
  \item Crowther, W.J.\ \& Potts, J.R.\ (2007). Simulation of a spin-stabilised sports disc. \emph{Sports Engineering} \textbf{10}(1), 3--21.
  \item Franke, J.\ et al.\ (2007). \emph{Best Practice Guideline for CFD Simulation of Flows in the Urban Environment.} COST Action 732.
  \item Hargreaves, D.M.\ \& Wright, N.G.\ (2007). On the use of the k-$\varepsilon$ model in commercial CFD software\ldots\ \emph{JWEIA} \textbf{95}, 355--369.
  \item Hubbard, M.\ \& Hummel, S.A.\ (2000). Simulation of Frisbee Flight. 5th Math \& Computers in Sport, Sydney.
  \item Hummel, S.A.\ (2003). \emph{Frisbee Flight Simulation and Throw Biomechanics.} MS Thesis, UC Davis.
  \item Kamaruddin, N.M., Potts, J.R.\ \& Crowther, W.J.\ (2018). Aerodynamic performance of flying discs. \emph{Aircraft Eng.\ Aerosp.\ Tech.}\ \textbf{90}(2), 390--397.
  \item Lissaman, P.B.S.\ \& Hubbard, M.\ (2010). Maximum range of flying discs. \emph{Procedia Engineering} \textbf{2}(2), 2529--2535.
  \item Lukes, R., Hart, J.\ et al.\ (2022). Disc golf trajectory modelling combining CFD and rigid body dynamics. \emph{Sports Engineering} \textbf{25}, 24.
  \item Potts, J.R.\ \& Crowther, W.J.\ (2002). Frisbee\texttrademark{} aerodynamics. AIAA 2002-3150.
  \item Richards, P.J.\ \& Hoxey, R.P.\ (1993). Appropriate boundary conditions for computational wind engineering models. \emph{JWEIA} \textbf{46--47}, 145--153.
  \item Tominaga, Y.\ et al.\ (2008). AIJ guidelines for practical applications of CFD to pedestrian wind environment. \emph{JWEIA} \textbf{96}, 1749--1761.
  \item Yang, Y.\ et al.\ (2009). New inflow boundary conditions for modelling the neutral equilibrium atmospheric boundary layer. \emph{JWEIA} \textbf{97}, 88--95.
  \item Shinozuka, M.\ \& Deodatis, G.\ (1991). Simulation of stochastic processes by spectral representation. \emph{Applied Mechanics Reviews} \textbf{44}(4), 191--204.
  \item Mann, J.\ (1998). Wind field simulation. \emph{Probabilistic Engineering Mechanics} \textbf{13}(4), 269--282.
  \item Garratt, J.R.\ (1990). The internal boundary layer --- a review. \emph{Boundary-Layer Meteorology} \textbf{50}, 171--203.
  \item OpenStreetMap contributors (2026). Building footprints, Loyola Lake Shore Campus. \url{https://www.openstreetmap.org} (ODbL).
  \item U.S.\ Geological Survey (2026). 3DEP Elevation Point Query Service. \url{https://epqs.nationalmap.gov}.
\end{description}

\end{document}
